\chapter{Experiments and Discussion}
\label{ch:experiments_discussion}

% ==============================================================================
% Core Questions:
% How do the proposed Predictors and MPC formulations perform in closed-loop against
% the ground-truth Plant, what are their operational trade-offs, and where do they break down?
% ==============================================================================

% Contents:
% - Experimental Protocol and Evaluation Criteria:
%   * In-silico patient scenarios: Spontaneous seizure generation in EZ and recruitment into PZ.
%   * Quantitative metrics: Seizure State, Seizure Burden, total injected charge, solve times.
%   * Uncontrolled baseline and reference conditions.
% - Study 0: Predictability experiment:
%   * what metrics can 1. detect seizures and 2. are predictable
%   * based of notebooks/state/scoring.py
%   * Note: the metric should be the MPC cost function? instead of what is currently there 
% - Study 1: Offline Predictor Benchmarking:
%   * Selection of STFT geometry: Trading temporal responsiveness against frequency resolution.
%   * Training loss comparison: Multi-step Rollout stability of spectral training loss vs. raw MSE loss.
%   * Prediction-target comparison: waveform Predictors trained with a spectral Loss vs. Predictors
%     that forecast spectral Observable Frames directly; compare Control Horizon accuracy and response to current.
%   * Model comparison: Linear DMDc vs. Nonlinear MLP over the Control Horizon.
% - Study 2: Numerical Solver and Transcription Benchmarking:
%   * Solve time, convergence reliability, and constraint satisfaction across horizon lengths.
%   * Single shooting vs. multiple shooting vs. NARX multiple shooting
% - Study 3: Closed-Loop Seizure Suppression:
%   * Proof of closed-loop efficacy: Attenuating seizure amplitude and preventing PZ recruitment via scalp EEG feedback.
%   * Waveform MPC vs. Spectral MPC: Empirical demonstration of why waveform tracking fails while spectral MPC succeeds.
%   * With the same spectral MPC Cost, compare waveform and direct spectral Predictors to isolate
%     the effect of the prediction target on closed-loop Seizure Suppression.
%   * Sparse Stimulation investigations: 
%     * compare L_1 loss, bonmin, block candidate, homotopy

% ---
% would be great if we could define some measure for "unnecessary stimulation" i.e. stimulation when it wasn't immenently necessary for suppression